\documentclass[10pt,a4paper]{amsart}
\usepackage[margin=19mm]{geometry}
\usepackage{amsmath,amssymb,mathtools}
\usepackage[scr=rsfs,cal=euler]{mathalfa}
\usepackage{xfrac}
\usepackage[unicode,hidelinks,bookmarks=false]{hyperref}
\newcommand{\ie}{i.e.\ }
\newcommand{\cf}{cf.\ }
\newcommand{\resp}{respectively\ }
\newcommand{\Subsection}{\subsection}
\providecommand{\nobreakdash}{\nobreak\hbox{-}}
\setlength{\emergencystretch}{2em}
\multlinegap0pt
\usepackage{bm}
\usepackage{xcolor}
\usepackage{tikz}
\usepackage{amssymb}
\usepackage{dsfont}
\usepackage[shortlabels]{enumitem}
\usepackage[normalem]{ulem}
\newtheorem{proposition}{Proposition}
\title{Some Taylor varieties with null Hessian}
\author{Thais Gomes Ribeiro, Elena Guardo, Manuela Muzika Dizdarevi\'c, Maryam Nowroozi, Pierpaola Santarsiero, Paola Supino}
\date{Source: arXiv:2605.03923v1 (CC BY 4.0)}
\begin{document}
\maketitle

\noindent\emph{Source and licence.} Some Taylor varieties with null Hessian. Version arXiv:2605.03923v1 (CC BY 4.0), distributed by arXiv under the Creative Commons Attribution 4.0 International licence, http://creativecommons.org/licenses/by/4.0/, as recorded in the arXiv licence metadata for this exact version and shown on the abstract page https://arxiv.org/abs/2605.03923v1. Full licence notice: \texttt{licenses/arxiv-2605.03923-CC-BY-4.0.txt}.\par\smallskip

\noindent\textit{Editorial scope.} Counted unit: the article's proposition nonsingular (source0.tex lines 415--417). The article's complete original proof of it is delivered with it (source0.tex lines 419--423, one \texttt{proof} environment of 5 lines). No mathematics is written, repaired or re-derived: the statement and the proof are the article's own lines, in the article's own notation, and no retained line is substituted or re-flowed.  No further source-local context is needed: every cross-reference of every retained line resolves inside the two delivered spans.  Nothing was omitted from any retained span, so the package carries no omission record.  No retained line contains a citation, so the wrapper carries no bibliography and no bibliography entry.  No retained line contains a figure environment, an image-inclusion command or drawing code, so no image is delivered and the package declares no asset dependency.  Editorial formatting only, in addition: an AMS \texttt{amsart} preamble that restates the article's own declarations and macros used by the retained lines, namely the environments \texttt{proposition} on separate counters, together with the standard packages those lines need.  The retained environments print in this document as Proposition 1 in order of appearance, whereas the article prints the same items with the numbering of its own full document.  The complete native source of the article as distributed in the arXiv e-print for this exact version is bundled unchanged as \texttt{sources/199830/source0.tex}.
\begin{proposition}\label{nonsingular}
Let $\mathcal{T}^{n}_{d,e,m} \subset \mathbb{P}^N$ be a Taylor variety. If $\mathcal{T}^{n}_{d,e,m}$ is a non-defective hypersurface, then the associated Padé matrix $P_T$ is non-singular, i.e., $\det(P_T) \not\equiv 0$.
\end{proposition}

\begin{proof}
By definition, a non-defective hypersurface in $\mathbb{P}^N$ has dimension $N-1$ and it is the zero set of a single irreducible polynomial $f \in \mathbb{C}[c_0, \dots, c_N]$. In the study of Taylor varieties, the conditions for a point to lie on $\mathcal{T}^{n}_{d,e,m}$ are given by the vanishing of the determinant of the Padé matrix $P_T$, which encodes the syzygies relations of the Taylor polynomial. 

If $P_T$ were singular for a general choice of coefficients, then $\det(P_T)$ would be the zero polynomial. This would imply that the variety has codimension greater than one, contradicting the hypothesis that $\mathcal{T}^{n}_{d,e,m}$ is a hypersurface. Therefore, the non-singularity of $P_T$ is a necessary condition for the variety to be a hypersurface whose defining equation is $f = \det(P_T)$.
\end{proof}

\end{document}
